% RQ4 Introduction: Treasury Resilience

\subsection{Motivation}

DAO treasuries represent significant pools of capital. As of 2023, the top 100 DAOs collectively held over \$25 billion in treasury assets \citep{deepdao2023analytics}. These treasuries fund operations, grants, development, and community initiatives---they are the financial foundation of decentralized organizations.

Yet treasury management in most DAOs remains primitive. Common patterns include:
\begin{itemize}
    \item \textbf{No reserves}: Treasury fully available for spending proposals
    \item \textbf{Concentration risk}: Treasury dominated by native governance token
    \item \textbf{Reactive management}: Policies developed only after crises
    \item \textbf{Informal limits}: Spending constrained by norms rather than rules
\end{itemize}

The consequences of poor treasury management have been visible: DAOs depleting reserves during bear markets, organizations unable to fund operations when token prices crashed, and treasuries drained by successive large proposals without regard for sustainability.

\subsection{Research Question}

This paper investigates:

\begin{quote}
\textbf{RQ4:} How do treasury management parameters (buffer fractions, emergency top-ups, spending limits) affect fiscal resilience?
\end{quote}

Specifically, we examine:
\begin{enumerate}
    \item How do buffer reserves affect treasury volatility and organizational sustainability?
    \item When should emergency top-up mechanisms activate, and how quickly should they restore reserves?
    \item What spending limits balance operational flexibility with fiscal prudence?
    \item How do these parameters interact under different market conditions?
\end{enumerate}

\subsection{Treasury Policy Dimensions}

We study three policy dimensions:

\subsubsection{Buffer Fractions}

The portion of treasury held as reserve, unavailable for proposal spending:
\begin{equation}
\text{available} = (1 - \text{buffer\_fraction}) \cdot \text{treasury}
\end{equation}

\subsubsection{Emergency Top-Ups}

Mechanisms that restore buffer when depleted below threshold:
\begin{itemize}
    \item Revenue diversion: Redirect incoming funds to buffer
    \item Spending freeze: Halt non-essential spending until restored
    \item External injection: Protocol fees or other sources
\end{itemize}

\subsubsection{Spending Limits}

Caps on outflows per period or per proposal:
\begin{equation}
\text{spend}(p) \leq \min(\text{requested}(p), \text{limit})
\end{equation}

\subsection{Contributions}

This paper contributes:

\begin{enumerate}
    \item \textbf{Resilience metrics} for treasury health assessment
    \item \textbf{Systematic analysis} of buffer, top-up, and limit effects
    \item \textbf{Optimal parameter ranges} for different risk profiles
    \item \textbf{Practical guidelines} for treasury policy design
\end{enumerate}

\subsection{Paper Organization}

Section~\ref{sec:background} reviews treasury management in DAOs. Section~\ref{sec:theory} develops formal models of treasury dynamics. Section~\ref{sec:architecture} describes simulation implementation. Section~\ref{sec:methodology} details experimental design. Section~\ref{sec:results} presents findings. Section~\ref{sec:discussion} interprets results for practitioners. Section~\ref{sec:conclusion} concludes.
